% !TeX root = ../papers/nc-p3.tex
\section{行列恒等式によるYoneda積の単射性と小指数の頂点拡張}
\label{sec:pi-yoneda-bridges}
\begin{lemma}\label{lem:pi-yoneda-bridges}
標準正則窓の各隅が $s$ 行列の恒等式を満たすとする。
互いに非同型な絶対単純 $S,T,L_1,\ldots,L_n$ を取る。
各 $L$ について、$S,T,L$ が全て非零となるある頂点でその三次元数の和が $s$ を超えるとする。
このときYoneda積
\[
 \bigoplus_L\operatorname{Ext}^1(S,L)\otimes\operatorname{Ext}^1(L,T)
                   \longrightarrow\operatorname{Ext}^2(S,T)
\]
は単射である。また $S,T$ だけの次元数の和がある頂点で $s$ を超えれば両向きの第一Extは零である。
\end{lemma}
\begin{proof}
零へ行くテンソルを取り、各成分のテンソル階数を $m_L$ とする。
$N=\bigoplus L^{m_L}$ を通る二拡張にまとめれば、積の消滅は一つの三段加群
（上から $S,N,T$）に組めることを意味する。短完全列のExt完全列を使えばよい。
この表現の有限次元像で半単純商の行列単位を持ち上げる。
各中間型で、$S$ から中間への根基の矩形行列は $m_L$ 個の重複方向を全て張る。
もし真部分空間ならそれを消す線型結合への押し出し拡張が分裂するが、
階数分解で採った $m_L$ 個の第一Extは独立である。半単純行列単位を左右から掛ければ
この重複空間と全矩形行列の積になる。中間から $T$ へも双対の全重複空間が出る。
従って $m_L>0$ なら指定頂点の隅で三つの行列ブロックを順に通る非零積を作れる。
頂点への切断は各ブロックをその頂点空間に圧縮するだけで、非零な三頂点空間の
矩形行列と行列単位を同じように選べる（頂点冪等元の分離代数を保って半単純持上げを取れる）。
総計 $s+1$ 個の直交対角単位を通る積があれば、単位と次への矢を交互に並べた
標準交代式 $s_{2s}$ に非零評価がある。定理\ref{thm:pi-periodic-point-paths}の行列単位の議論による。
矛盾なので全成分が零である。二段非分裂拡張にも同じ証明を使う。
\end{proof}

\begin{theorem}\label{thm:one-sided-balanced-small-simples}
四辺有限の正則窓で各隅が四行列の恒等式を満たすとする。頂点次元が全て $d=2$ または全て $d=3$ の
任意の単純 $S$ は、全頂点への第一Extが零か、全頂点からの第一Extが零である。
一頂点での両方向の同時非消滅を排除するだけでなく、異なる頂点に出入りする場合も排除する。
また次元が2と3に分かれる単純のうち $(2,2,3,3),(2,3,2,3),(2,3,3,3)$ およびその巡回移動は現れない。
\end{theorem}
\begin{proof}
$S_\lambda$ を一周倍率 $\lambda$ の移動とする。同型倍率は有限群である。
ある正長語のtraceは非零であり、同型ならその文字長の倍率冪が1だからである。
一般の異なる二倍率では $\operatorname{Ext}^1(S_\lambda,S_\mu)=0$ である。
実際もし一般の比で次元が正なら、固定した一般の二端の間に任意個の異なる倍率を挿入でき、
両側のExtはともに正次元となる。Ext次元は有限Koszul複体の階数による構成可能関数である。
$3d>4$ なので補題の単射性は有限次元の第二Extに任意個の独立な元を要求する。
一般の二端で逆向きも零となり、四次元CY双対と定数次元ベクトルのEuler数零から
$\operatorname{Ext}^2(S,S_\lambda)=0$ も一般倍率で得る。

$a_i=\dim\operatorname{Ext}^1(S,k_i)$、$b_i=\dim\operatorname{Ext}^1(k_i,S)$ とする。
定数次元の場合、両向きの第二Extの次元は $a_i+b_i$ である。
補題を一般の $S,T=S_\lambda$ と中間 $k_i$ に使い、$2d+1>4$ より $a_i b_i=0$。
このとき $a_i,b_i\le1$ ($d=2$)、または $\le3$ ($d=3$) である。
例えば行拡張は入射辺での行解空間から元の $d$ 行を除いたものである。
行解空間は $4$ 行 $d$ 列の階数一を含めない。含めば前の矢を一共変元で評価した像が二次元以上あるため
辺グラフに直線が出る。非薄単純での全ベクトル・双対共変元の評価次元下限2を用いた。
従って行解空間はSegreの余次元 $3(d-1)$ 以下で、列も同じである。

両方向が別々に存在するとし、$a_i>0,b_j>0$ の間の最短順距離を $l$ とする。
右表現で矢を順に作用させる。$1\le l\le3$、途中の全ての $a,b$ は零である。
$i$ の頂点射影の長さ $l$ 以上を切り落とした加群を $P$ とする。頭は $k_i$、
その後ろは途中の頂点単純だけである。$W=(e_i B_l e_j)^*$ の任意の $w$ は
最後の一段を $k_j$ に評価する延長 $0\to k_j\to N_w\to P\to0$ を定める。
その類 $\eta_w$ は $w$ に線型で、$w\ne0$ なら頭から末尾への長さ $l$ の作用が非零である。
途中の第二Ext消滅により $\operatorname{Ext}^1(S,P)\to\operatorname{Ext}^1(S,k_i)$ は全射なので
その線型切断を選ぶ。また $\dim\operatorname{Ext}^2(P,T)\le a_i$。
$S$ から $N_w$ への拡張を作る障害と、$N_w$ から $T$ への拡張の障害はそれぞれ
この切断と $\eta_w$ の積、$\eta_w$ と $\operatorname{Ext}^1(k_j,T)$ の積である。
従って $w$ に線型な、大きさ $b_j\times a_i$ と高々 $a_i\times b_j$ の二行列である。

両方に非零核を与える $w\ne0$ がある。$a_i=b_j$ なら二行列式の共通零点を取り、
足りない行は零で補ってよい。$a_i\ne b_j$ なら広い方は常に核があり、他方の非単射条件の
行列空間での余次元は $|a_i-b_j|+1\le3$。$\dim W=\binom{l+3}3\ge4$ なので
写像の核があればそれを、なければ射影次元定理を使う。こうして両端に非分裂な第一拡張を保った
二つの延長を得る。最後に $\operatorname{Ext}^2(S,T)=0$ によって三段 $S,N_w,T$ に組む。
中間の頭 $k_i$ と末尾 $k_j$ はこの表現に各一度しか現れず、その間の経路は非零である。
有限像の半単純ブロックを切り出せば、最初と最後の非分裂矩形とこの経路を連接できる。
とくに $i$ 番隅で $S$ の $d$ 個、$k_i$ の一個、$T$ の $d$ 個の直交単位を非零積で通れる。
途中で $j$ 番を経る後半は全体を一つの矩形元とし、両端を $i$ 番で圧縮すればよい。
これは再び四行列の恒等式に反する。

最後に混合次元で $d_i=\dim S_i$ とする。$a_i\le3(d_{i-1}-1)-d_i$、
$b_i\le3(d_{i+1}-1)-d_i$ は同じ行・列計算による。
$h_i=2d_i-4d_{i-1}-4d_{i+1}+6d_{i+2}$ とすると $a_i+b_i+h_i\ge0$。
異なる倍率間の第一Extは一つの三次元頂点だけで $3+3>4$ となり消えるので、
第二Extは $\sum d_i h_i$ 次元である。従って補題から $\sum a_i b_i\le\sum d_i h_i$。
$(2,2,3,3)$ では右辺が $-4$、$(2,3,2,3)$ では右辺16に対し二つの二次元位置で
$a_i=b_i=4$、$(2,3,3,3)$ では右辺2に対し中央の三次元位置で $a_i=b_i=3$ が必要となり、
いずれも不可能である。混合型で残る可能性は $(2,2,2,3)$ の巡回移動だけである。
\end{proof}

\begin{corollary}\label{cor:semiprime-three-matrix-clifford}
全四隅が三行列の恒等式を満たす有限四辺の標準正則窓について、一隅が半素であれば
必ず通常のClifford窓である。これは十五次元閉包の仮定を使わない。
\end{corollary}
\begin{proof}
折り畳みもPIである。任意の非頂点単純は薄いか全頂点が次元2以上であり、
どこかで3に達すれば非零作用する正次数の三行列中心多項式がその隅から全中心へ延びる。
従って頂点と全Extが消えEuler式から $3\delta$ となる。それ以外の非薄は $2\delta$ である。
薄い単純の頂点との拡張は両向きに零である。

各斉次極小素商 $R$ のtrace環 $T=RC$ を取る。$C$ は有限型可換、$T$ はその上有限である。
非原点の $C$ の閉指標で $T$ の繊維単純が頂点単純を含まないことを示す。
一般の絶対単純からの特殊化半単純型には繊維の全単純型が出る
（定理\ref{thm:pi-vertex-extension-criterion}のCayley--Hamiltonの議論）。
一般の次元は高々 $3\delta$ なので、頂点と非頂点をともに含むとき、
頂点と拡張を持ち得る非薄はただ一つの $2\delta$ だけである。
前定理からその頂点への、または頂点からの全第一Extが零であり、
薄い因子も含めた二群の間で少なくとも一向きの第一Extが全て零になる。
この半単純型は絶対単純に平滑化できない。実際表現多様体で半単純軌道は一般単純の閉包にあり
（特殊格子の軌道から半単純化へ退化してもよい）、そこへ曲線で近づく。
完備離散付値環の各厚みで消えている向きの非対角ブロックを消せる。
最初に残るその係数は二群間の一次コサイクルで、基底変更で除ける。反対側のブロックがあっても
下三角ブロックの方程式には、既に消した向きなしの項は生じない。
極限で真の不変部分ができ矛盾する。非原点で全因子が頂点単純でもtraceが全て零となるため不可能である。

従って $T/T R_+$ は $C$ の原点だけに支えられ有限次元である。
非原点で非零ならその有限次元繊維に、$R_+$ が消す単純右商があり、繊維上は $C$ がスカラーなので
これは上で除いた頂点型 $T$-単純となる。ここで $TR_+$ は右 $R$ 加群かつ $C$ 加群であり、
$T=RC$ だから右 $T$ 加群をなす。次数帰納で $T$ は $R$ 上有限生成となり、
$T=RC$ のため有限個の中心係数を加群生成元にできる。
有限中心化拡大のNoether降下定理（Formanek--Jategaonkar）により、
Noether環 $T$ のこの部分環 $R$ は左右Noetherである。零次元素商も自明である。

一隅半素なら全体も半素である。非零冪零両側イデアルがあれば、その非零元を左右から矢で
指定頂点まで伸ばせる。途中で全矢に消される元がないことは左右ASの頂点socle消滅による。
結果は指定隅の非零冪零イデアルとなって矛盾する。半素アフィンPI環は有限個の斉次極小素商の
積に埋め込まれる。これらは全て左右の $B$ 加群としてNoetherなので $B$ 自体も左右Noetherである。
よって中心有限かつ合成整域で、頂点指数3は不可能、指数1は有限四辺で不可能、
指数2は既出の通常Clifford分類で尽くされる。
\end{proof}
